
\subsubsection{4.1 Research Questions}

The study was designed to answer the following pre-registered research questions:

\begin{itemize}
\item \textbf{RQ1 (Existence):} Does prescreening APPS introductory problems yield $\geq 50$ qualifying problems with S\_term $\in$ [0.3, 0.55]?
\item \textbf{RQ2 (Variance Gate):} In qualifying groups, does E[Var(r\_ratio)] / E[Var(r\_binary)] $\geq$ 1.$5\times$?
\item \textbf{RQ3 (Advantage Diversity):} Does R\_ratio produce $\geq 5$ distinct advantage levels per group versus ~2 under R\_binary?
\item \textbf{RQ4 (ZRF Escape):} Does R\_ratio produce earlier escape from the Zero Reward Fraction plateau, specifically ZRF\_ratio(t\textit{) < 0.8 $\times$ ZRF\_binary(t}) with log-rank p < 0.05?
\end{itemize}

RQ1 and RQ2 correspond to Stage 1 (h-e1). RQ3 and RQ4 correspond to Stage 2 (h-m1 through h-m4).

\subsubsection{4.2 Dataset}

The APPS benchmark \citep{Hendrycks2021} was used, restricted to the introductory difficulty tier (difficulty=0) with a minimum test count filter of $T\geq 3$. This yields 1,923 qualifying problems from the full dataset of 2,639 training problems. A random sample of 300 problems (seed=42) was processed for prescreening, representing 15.6\% of available introductory problems.

\subsubsection{4.3 Model and Baselines}

\textbf{Model:} Qwen2.5-Coder-7B-Instruct \citep{Hui2024}, using the base instruction-tuned checkpoint without any APPS-specific fine-tuning. An SFT checkpoint fine-tuned on APPS introductory solutions was specified as a controlled variable in the hypothesis design but was not available at experiment time.

\textbf{Primary baseline:} R\_binary --- standard GRPO binary execution reward.

\textbf{External reference:} Afterburner \citep{Du2025} serves as a consistency reference for the SFT prerequisite finding.

\subsubsection{4.4 Implementation Details}

\textbf{Code modules:} prescreening.py (371 lines), evaluate.py, reward\_fn.py, data\_loader.py, execution\_sandbox.py, visualization.py. All modules implemented in Python.

\textbf{Execution sandbox:} Isolated subprocess execution with a 5-second per-test-case timeout.

\textbf{Sampling parameters:} k=8 rollouts per problem, temperature $\tau =0.8$, max\_new\_tokens=1,024, seed=42.

\textbf{Training framework (planned for Stage 2):} HuggingFace TRL GRPOTrainer v0.29.0. GRPO algorithm modifications were not planned; only the reward function was varied between conditions.

\textbf{Hardware:} Single NVIDIA H100 NVL GPU (80 GB). Stage 1 prescreening of 300 problems with k=8 rollouts completed in approximately 42 minutes.

\subsubsection{4.5 Gate Metrics (Pre-registered)}

Three gate metrics were pre-specified for the Stage 1 gate (h-e1 MUST\_WORK gate):

\begin{itemize}
\item fraction\_k\_pass\_ge1 $\geq$ 0.10: fraction of processed problems with at least one rollout passing at least one test case.
\item pct\_groups\_above\_1.5x $\geq$ 0.80: fraction of qualifying groups where Var(r\_ratio) / Var(r\_binary) $\geq$ 1.5.
\item n\_prescreened $\geq$ 50: count of problems with S\_term $\in$ [0.3, 0.55].
\end{itemize}

All three thresholds were specified before running experiments.

